% ===========================================================================
%  Chapitre 4 — Limites et branches infinies
%  PE 2026, composante ANALYSE
% ===========================================================================
\bmtchapitre{Limites et branches infinies}
  {Calculer la limite d'une fonction aux bornes de son ensemble de
   définition, lever une forme indéterminée, et interpréter graphiquement le
   résultat : asymptote verticale, horizontale ou oblique, branche
   parabolique.}
  {Analyse \textbullet\ environ 5 heures}

Une limite répond à une question de comportement : que devient $f(x)$ quand
$x$ s'approche d'une valeur interdite, ou quand $x$ devient très grand ? La
réponse ne s'invente pas, elle se calcule — et surtout, elle se \emph{lit
ensuite sur le dessin}. C'est là tout l'intérêt du chapitre : chaque limite
calculée devient une droite tracée, asymptote de la courbe.

À l'examen, le problème commence presque toujours par « calculer les limites
de $f$ aux bornes de son ensemble de définition », et cette question vaut
un point ou deux, avant même d'avoir dérivé quoi que ce soit. Elle se prépare
donc à part.

% ---------------------------------------------------------------------------
\begin{revision}
Aucune notion de ce chapitre n'est utilisée ici.

\begin{enumerate}[label=\textbf{\arabic*.}, leftmargin=*, itemsep=0.35em]
  \item Pour quelles valeurs de $x$ le quotient $\dfrac{2x+1}{x-3}$
        existe-t-il ?
  \item Factoriser $x^{2}-5x+6$.
  \item Simplifier $\dfrac{x^{2}-4}{x+2}$ pour $x \neq -2$.
  \item Effectuer $\dfrac{3x^{2}-x}{x^{2}}$ et simplifier.
  \item Calculer $\dfrac{1}{0{,}001}$ et $\dfrac{1}{1\,000\,000}$.
  \item Que vaut $\dfrac{5}{x}$ quand $x$ est très grand ? très proche de
        $0$ ?
\end{enumerate}
\end{revision}

\begin{automatismes}
Sans calculatrice.

\begin{enumerate}[label=\textbf{\alph*)}, leftmargin=*, itemsep=0.25em]
  \item $\dfrac{1}{x}$ quand $x \to +\infty$
  \item $\dfrac{1}{x^{2}}$ quand $x \to +\infty$
  \item $x^{2}$ quand $x \to -\infty$
  \item $x^{3}$ quand $x \to -\infty$
  \item $\dfrac{3x}{x}$ pour $x \neq 0$
  \item $\dfrac{2x^{2}}{x^{2}}$ pour $x \neq 0$
  \item Le signe de $x-3$ quand $x$ est un peu plus petit que $3$
  \item Le signe de $x-3$ quand $x$ est un peu plus grand que $3$
\end{enumerate}
\end{automatismes}

\begin{corrige}[automatismes]
a) tend vers $0$ \quad b) tend vers $0$ \quad c) tend vers $+\infty$ \quad
d) tend vers $-\infty$ \quad e) $3$ \quad f) $2$ \quad g) négatif \quad
h) positif.
\end{corrige}

% ===========================================================================
\section{Limite en l'infini, limite en un point}

\begin{definition}[limite en $+\infty$]
Dire que $f$ a pour limite le réel $\ell$ en $+\infty$, ce qui s'écrit
$\limpinf f(x) = \ell$, signifie que $f(x)$ devient aussi proche que l'on
veut de $\ell$ dès que $x$ est assez grand. Dire que
$\limpinf f(x) = +\infty$ signifie que $f(x)$ dépasse tout nombre fixé à
l'avance dès que $x$ est assez grand.
\end{definition}

\begin{definition}[limite en un réel]
Soit $a$ un réel. On note $x \to a^{+}$ lorsque $x$ tend vers $a$ en restant
\emph{supérieur} à $a$, et $x \to a^{-}$ lorsque $x$ tend vers $a$ en restant
\emph{inférieur} à $a$. Ces deux limites peuvent être différentes, et c'est
même le cas le plus fréquent aux bornes interdites.
\end{definition}

\begin{propriete}[limites à connaître par c\oe ur]
\[
  \limpinf \frac{1}{x} = 0, \qquad
  \limminf \frac{1}{x} = 0, \qquad
  \lim_{x \to 0^{+}} \frac{1}{x} = +\infty, \qquad
  \lim_{x \to 0^{-}} \frac{1}{x} = -\infty .
\]
Pour tout entier $n \geq 1$ : $\limpinf x^{n} = +\infty$, et
$\limminf x^{n} = +\infty$ si $n$ est pair, $-\infty$ si $n$ est impair.
\end{propriete}

\begin{visualisation}[la même fonction, quatre limites]
\begin{center}
\begin{tikzpicture}
\begin{axis}[
  width = 11.6cm, height = 6.4cm,
  axis lines = middle, xlabel = {$x$}, ylabel = {$y$},
  xmin = -4.6, xmax = 6.6, ymin = -4.4, ymax = 5.4,
  xtick = {-4,-2,2,4,6}, ytick = {-4,-2,2,4},
  axis line style = {bmtGris},
  tick label style = {font=\scriptsize, color=bmtGris},
  clip = true, restrict y to domain = -8:8,
]
  \addplot[bmtIndigo, line width=1.4pt, domain=-4.5:0.94, samples=200]
    {1+2/(x-1)};
  \addplot[bmtIndigo, line width=1.4pt, domain=1.06:6.5, samples=200]
    {1+2/(x-1)};
  \addplot[bmtLaterite, line width=1pt, dash pattern=on 3pt off 2pt,
           domain=-4.5:6.5] {1};
  \draw[bmtLaterite, line width=1pt, dash pattern=on 3pt off 2pt]
    (axis cs:1,-4.3) -- (axis cs:1,5.3);
  \node[bmtLaterite, font=\scriptsize, anchor=south west]
    at (axis cs:4.4,1.08) {$y = 1$};
  \node[bmtLaterite, font=\scriptsize, anchor=south west]
    at (axis cs:1.12,3.7) {$x = 1$};
\end{axis}
\end{tikzpicture}
\end{center}

La courbe de $f(x) = 1+\dfrac{2}{x-1}$ raconte quatre limites d'un coup :
quand $x \to +\infty$ et quand $x \to -\infty$, la courbe se colle à la
droite horizontale $y = 1$ ; quand $x \to 1^{+}$, elle file vers le haut ;
quand $x \to 1^{-}$, vers le bas. Les deux droites en pointillé ne font pas
partie de la courbe : ce sont ses \emph{asymptotes}.
\end{visualisation}

\begin{propriete}[opérations sur les limites]
Les limites d'une somme, d'un produit, d'un quotient se calculent en
remplaçant chaque terme par sa limite, avec les règles habituelles du calcul
sur $\pm\infty$ — sauf dans quatre cas, dits \textbf{formes indéterminées} :
\[
  \infty-\infty, \qquad 0\times\infty, \qquad \frac{\infty}{\infty},
  \qquad \frac{0}{0}.
\]
\end{propriete}

\begin{avertissement}
« Forme indéterminée » ne veut pas dire « pas de limite » : cela veut dire
que le calcul direct ne permet pas de conclure et qu'il faut transformer
l'écriture. Écrire « F.I. » sur une copie et s'arrêter là ne rapporte aucun
point.
\end{avertissement}

% ===========================================================================
\section{Lever une indétermination}

\begin{methode}{le terme de plus haut degré, à l'infini}
Pour une fonction polynôme ou rationnelle, quand $x \to \pm\infty$ :
\begin{enumerate}[label=\textbf{\arabic*.}, leftmargin=*, itemsep=0.15em]
  \item Mettre en facteur le terme de plus haut degré au numérateur et au
        dénominateur.
  \item Simplifier.
  \item Conclure : tous les termes en $\frac{1}{x^{k}}$ tendent vers $0$.
\end{enumerate}
On retient le résultat : \emph{à l'infini, un polynôme se comporte comme son
terme de plus haut degré}, et un quotient de polynômes comme le quotient de
leurs termes de plus haut degré.
\end{methode}

\begin{exemple}[trois quotients, trois issues]
\textbf{a)} $\limpinf \dfrac{3x^{2}-x+1}{x^{2}+5}
= \limpinf \dfrac{x^{2}\left(3-\frac1x+\frac{1}{x^{2}}\right)}
{x^{2}\left(1+\frac{5}{x^{2}}\right)}
= \dfrac{3}{1} = 3$ — le degré est le même en haut et en bas.

\textbf{b)} $\limpinf \dfrac{2x+7}{x^{2}-1} = \limpinf \dfrac{2}{x} = 0$ —
le bas l'emporte.

\textbf{c)} $\limpinf \dfrac{x^{3}+1}{4x-2} = \limpinf \dfrac{x^{2}}{4}
= +\infty$ — le haut l'emporte.
\end{exemple}

\begin{methode}{la forme $\frac{0}{0}$ en un point}
Si le numérateur et le dénominateur s'annulent tous les deux en $a$, c'est
que $(x-a)$ les divise tous les deux : on factorise, on simplifie, puis on
remplace.
\end{methode}

\begin{exemple}[factoriser pour simplifier]
$\displaystyle\lim_{x \to 2} \frac{x^{2}-4}{x^{2}-5x+6}$ donne, par
remplacement direct, $\frac00$. Or
$x^{2}-4 = (x-2)(x+2)$ et $x^{2}-5x+6 = (x-2)(x-3)$, donc pour $x \neq 2$,
\[
  \frac{x^{2}-4}{x^{2}-5x+6} = \frac{x+2}{x-3}
  \quad \xrightarrow[\;x \to 2\;]{} \quad \frac{4}{-1} = -4 .
\]
\end{exemple}

\begin{methode}{la limite en une valeur interdite}
Quand le dénominateur tend vers $0$ sans que le numérateur s'annule, la
limite est infinie — reste à trouver le signe. On dresse une petite ligne de
signes du dénominateur autour de la valeur interdite, et l'on combine avec
le signe du numérateur.
\end{methode}

\begin{exemple}[la ligne de signes qui décide]
Cherchons les limites de $f(x) = \dfrac{x+5}{x-2}$ en $2$.

Le numérateur tend vers $7 > 0$. Le dénominateur tend vers $0$ : il est
négatif à gauche de $2$, positif à droite. Donc
\[
  \lim_{x \to 2^{-}} f(x) = -\infty, \qquad
  \lim_{x \to 2^{+}} f(x) = +\infty .
\]
\end{exemple}

\begin{erreurclassique}
Une limite « en $2$ » n'existe pas toujours : il faut préciser $2^{-}$ ou
$2^{+}$ dès que la fonction change de comportement de part et d'autre.
Écrire $\displaystyle\lim_{x\to2} \frac{x+5}{x-2} = \infty$ sans signe ni
côté vaut zéro point.
\end{erreurclassique}

\begin{entrainement}[calculs de limites]
\exo[\niveau{1}] $\limpinf \left(x^{2}-3x+1\right)$

\exo[\niveau{1}] $\limminf \left(-2x^{3}+x\right)$

\exo[\niveau{1}] $\limpinf \dfrac{5x-1}{2x+3}$

\exo[\niveau{2}] $\limpinf \dfrac{x+4}{x^{2}+1}$

\exo[\niveau{2}] $\limminf \dfrac{2x^{2}-1}{x+1}$

\exo[\niveau{2}] $\displaystyle\lim_{x \to 3^{+}} \dfrac{1}{x-3}$ puis
$\displaystyle\lim_{x \to 3^{-}} \dfrac{1}{x-3}$

\exo[\niveau{2}] $\displaystyle\lim_{x \to 1} \dfrac{x^{2}-1}{x-1}$

\exo[\niveau{3}] $\displaystyle\lim_{x \to -2} \dfrac{x^{2}+3x+2}{x^{2}-4}$

\exo[\niveau{3}] $\displaystyle\lim_{x \to 0^{+}}
\left(\dfrac{1}{x}-\dfrac{1}{x^{2}}\right)$
\end{entrainement}

\begin{corrige}[calculs de limites]
\textbf{1.} $+\infty$ (comportement de $x^{2}$).
\textbf{2.} $+\infty$ : $-2x^{3} \to +\infty$ quand $x \to -\infty$.
\textbf{3.} $\frac52$.
\textbf{4.} $0$.
\textbf{5.} $\limminf \frac{2x^{2}}{x} = \limminf 2x = -\infty$.
\textbf{6.} $+\infty$ puis $-\infty$.
\textbf{7.} $\frac{(x-1)(x+1)}{x-1} = x+1 \to 2$.
\textbf{8.} $\frac{(x+1)(x+2)}{(x-2)(x+2)} = \frac{x+1}{x-2} \to
\frac{-1}{-4} = \frac14$.
\textbf{9.} $\frac{1}{x}-\frac{1}{x^{2}} = \frac{x-1}{x^{2}}$ ; le
numérateur tend vers $-1$, le dénominateur vers $0^{+}$ : la limite est
$-\infty$.
\end{corrige}

% ===========================================================================
\section{Interprétation graphique : les asymptotes}

\begin{definition}[asymptote verticale]
Si $\displaystyle\lim_{x \to a^{+}} f(x) = \pm\infty$ ou
$\displaystyle\lim_{x \to a^{-}} f(x) = \pm\infty$, la droite d'équation
$x = a$ est \textbf{asymptote verticale} à $\Cf$.
\end{definition}

\begin{definition}[asymptote horizontale]
Si $\limpinf f(x) = \ell$ (ou $\limminf f(x) = \ell$), la droite d'équation
$y = \ell$ est \textbf{asymptote horizontale} à $\Cf$ en $+\infty$
(respectivement en $-\infty$).
\end{definition}

\begin{definition}[asymptote oblique]
S'il existe deux réels $a \neq 0$ et $b$ tels que
\[
  \limpinf \left[f(x)-(ax+b)\right] = 0,
\]
la droite d'équation $y = ax+b$ est \textbf{asymptote oblique} à $\Cf$ en
$+\infty$. Même définition en $-\infty$.
\end{definition}

\begin{methode}{trouver une asymptote oblique}
Deux situations, deux gestes :
\begin{itemize}[leftmargin=1.3em, itemsep=0.2em]
  \item \emph{L'énoncé donne la droite} : on calcule
        $f(x)-(ax+b)$, on montre que cette différence tend vers $0$, et
        c'est fini. C'est la question la plus fréquente à l'examen.
  \item \emph{L'énoncé ne la donne pas} : pour une fonction rationnelle dont
        le degré du numérateur dépasse d'exactement un celui du
        dénominateur, on effectue la division et l'on écrit
        $f(x) = ax+b+\dfrac{\text{reste}}{\text{dénominateur}}$ ; le dernier
        terme tend vers $0$, les deux premiers donnent l'asymptote.
\end{itemize}
\end{methode}

\begin{exemple}[une asymptote oblique par division]
Soit $f(x) = \dfrac{x^{2}-3x+4}{x-2}$ sur $\R\setminus\{2\}$. On vérifie que
\[
  f(x) = x-1+\frac{2}{x-2}.
\]
En effet, $(x-1)(x-2)+2 = x^{2}-3x+2+2 = x^{2}-3x+4$. \checkmark

Comme $\dfrac{2}{x-2} \to 0$ en $\pm\infty$, la droite $(D) : y = x-1$ est
asymptote oblique à la courbe en $+\infty$ et en $-\infty$. De plus, le
signe de $\frac{2}{x-2}$ donne la position : la courbe est au-dessus de
$(D)$ pour $x > 2$, en dessous pour $x < 2$.
\end{exemple}

\begin{definition}[branche parabolique]
Si $\limpinf f(x) = \pm\infty$ et $\limpinf \dfrac{f(x)}{x} = \pm\infty$, on
dit que $\Cf$ admet une \textbf{branche parabolique de direction} $(Oy)$.
Si $\limpinf \dfrac{f(x)}{x} = 0$, la branche parabolique est de direction
$(Ox)$.
\end{definition}

\begin{visualisation}[les trois asymptotes et la branche parabolique]
\begin{center}
\begin{tikzpicture}
\begin{axis}[
  name = A, width = 4.5cm, height = 4.1cm, axis lines = middle,
  xmin = -0.4, xmax = 4.4, ymin = -0.4, ymax = 4.4,
  xtick = \empty, ytick = \empty, axis line style = {bmtGris},
  title = {\scriptsize horizontale}, title style = {color=bmtGris},
  clip = true, restrict y to domain = -1:6,
]
  \addplot[bmtIndigo, line width=1.2pt, domain=0.35:4.3, samples=120]
    {1.2+1.1/x};
  \addplot[bmtLaterite, line width=0.9pt, dash pattern=on 2.5pt off 2pt,
           domain=0:4.3] {1.2};
\end{axis}
\begin{axis}[
  name = B, at = {($(A.east)+(0.75cm,0)$)}, anchor = west,
  width = 4.5cm, height = 4.1cm, axis lines = middle,
  xmin = -0.4, xmax = 4.4, ymin = -0.4, ymax = 4.4,
  xtick = \empty, ytick = \empty, axis line style = {bmtGris},
  title = {\scriptsize verticale}, title style = {color=bmtGris},
  clip = true, restrict y to domain = -1:6,
]
  \addplot[bmtIndigo, line width=1.2pt, domain=1.62:4.3, samples=140]
    {1.1/(x-1.5)};
  \draw[bmtLaterite, line width=0.9pt, dash pattern=on 2.5pt off 2pt]
    (axis cs:1.5,-0.35) -- (axis cs:1.5,4.3);
\end{axis}
\begin{axis}[
  name = C, at = {($(B.east)+(0.75cm,0)$)}, anchor = west,
  width = 4.5cm, height = 4.1cm, axis lines = middle,
  xmin = -0.4, xmax = 4.4, ymin = -0.4, ymax = 4.4,
  xtick = \empty, ytick = \empty, axis line style = {bmtGris},
  title = {\scriptsize oblique}, title style = {color=bmtGris},
  clip = true, restrict y to domain = -1:6,
]
  \addplot[bmtIndigo, line width=1.2pt, domain=0.5:4.3, samples=140]
    {0.85*x+1/x};
  \addplot[bmtLaterite, line width=0.9pt, dash pattern=on 2.5pt off 2pt,
           domain=0:4.3] {0.85*x};
\end{axis}
\end{tikzpicture}
\end{center}

Trois façons pour une courbe de « suivre » une droite au loin. Dans le
quatrième cas — la branche parabolique — la courbe part vers l'infini sans
suivre aucune droite : elle s'écarte de toutes.
\end{visualisation}

\begin{entrainement}[asymptotes]
\exo[\niveau{1}] $f(x) = \dfrac{3x-1}{x+2}$ : donner les deux asymptotes.

\exo[\niveau{2}] $g(x) = 2+\dfrac{1}{x-4}$ : donner les asymptotes et la
position de la courbe par rapport à l'asymptote horizontale.

\exo[\niveau{2}] $h(x) = \dfrac{x^{2}+1}{x}$ : vérifier que
$h(x) = x+\dfrac1x$, en déduire une asymptote oblique et une asymptote
verticale.

\exo[\niveau{3}] $k(x) = \dfrac{2x^{2}-x+1}{x-1}$ : montrer que
$k(x) = 2x+1+\dfrac{2}{x-1}$, puis donner toutes les asymptotes.

\exo[\niveau{3}] $p(x) = x^{2}-3x$ : montrer que $\Cf$ admet une branche
parabolique de direction $(Oy)$ en $+\infty$.
\end{entrainement}

\begin{corrige}[asymptotes]
\textbf{1.} $x = -2$ (verticale) et $y = 3$ (horizontale).
\textbf{2.} $x = 4$ et $y = 2$. Le terme $\frac{1}{x-4}$ est positif pour
$x>4$ : la courbe est au-dessus de $y = 2$ à droite de $4$, en dessous à
gauche.
\textbf{3.} $\frac{x^{2}+1}{x} = \frac{x^{2}}{x}+\frac1x = x+\frac1x$ ;
comme $\frac1x \to 0$, la droite $y = x$ est asymptote oblique ; et $x = 0$
est asymptote verticale.
\textbf{4.} $(2x+1)(x-1)+2 = 2x^{2}-x-1+2 = 2x^{2}-x+1$ \checkmark ; d'où
$y = 2x+1$ asymptote oblique et $x = 1$ asymptote verticale.
\textbf{5.} $p(x) \to +\infty$ et $\frac{p(x)}{x} = x-3 \to +\infty$ : la
branche est parabolique de direction $(Oy)$.
\end{corrige}

\begin{qcm}
\exo $\limpinf \dfrac{4x^{2}+1}{2x^{2}-x}$ vaut :
\textbf{a)} $0$ \quad \textbf{b)} $2$ \quad \textbf{c)} $+\infty$ \quad
\textbf{d)} $4$ \reponseqcm{b}

\exo Si $\limminf f(x) = 3$, la courbe admet :
\textbf{a)} l'asymptote verticale $x = 3$ \quad
\textbf{b)} l'asymptote horizontale $y = 3$ en $-\infty$ \quad
\textbf{c)} une branche parabolique \quad \textbf{d)} rien de particulier
\reponseqcm{b}

\exo $\displaystyle\lim_{x \to 1^{-}} \dfrac{2}{x-1}$ vaut :
\textbf{a)} $+\infty$ \quad \textbf{b)} $-\infty$ \quad \textbf{c)} $2$
\quad \textbf{d)} $0$ \reponseqcm{b}

\exo Laquelle de ces écritures est une forme indéterminée ?
\textbf{a)} $\frac{5}{0^{+}}$ \quad \textbf{b)} $\infty+\infty$ \quad
\textbf{c)} $\frac{\infty}{\infty}$ \quad \textbf{d)} $0\times5$
\reponseqcm{c}

\exo Si $f(x) = 3x-2+\dfrac{1}{x}$, la droite $y = 3x-2$ est :
\textbf{a)} asymptote verticale \quad \textbf{b)} asymptote horizontale
\quad \textbf{c)} asymptote oblique \quad \textbf{d)} tangente à la courbe
\reponseqcm{c}
\end{qcm}

% ===========================================================================
\begin{annales}
Le problème de la série~A repose, à chaque session, sur le logarithme ou
l'exponentielle : ses questions de limites ne peuvent donc pas figurer ici
sans enfreindre l'ordre du livre. Les énoncés ci-dessous sont des exercices
réels portant sur des fonctions rationnelles, tirés de sujets africains de
niveau première et d'un baccalauréat blanc malgache, chacun avec sa
provenance exacte. Vous retrouverez les limites des sujets malgaches à leur
place, aux chapitres 7 et 8.

\exoann{Baccalauréat probatoire, Cameroun, séries D et TI, session 2023 —
exercice 4 (extrait)}
Soit $f$ la fonction définie sur $\R\setminus\{-1\}$ par
$f(x) = x+\dfrac{1}{x+1}$.
\begin{enumerate}[label=\textbf{\alph*)}, leftmargin=*, itemsep=0.15em]
  \item Calculer les limites de $f$ en $-\infty$ et en $+\infty$.
  \item Calculer les limites de $f$ en $-1^{-}$ et en $-1^{+}$.
  \item En déduire que la courbe admet deux asymptotes, dont on donnera les
        équations.
\end{enumerate}

\exoann{Baccalauréat probatoire, Cameroun, série D, session 2022 —
exercice 4 (extrait)}
Soit $f$ définie par $f(x) = \dfrac{3x}{3+4x}$.
\begin{enumerate}[label=\textbf{\alph*)}, leftmargin=*, itemsep=0.15em]
  \item Déterminer l'ensemble de définition de $f$.
  \item Calculer $\limpinf f(x)$ et $\limminf f(x)$, puis en déduire une
        asymptote.
  \item Étudier les limites à gauche et à droite de la valeur interdite.
\end{enumerate}

\exoann{Baccalauréat blanc, Lycée Philibert Tsiranana, session 2022-2023 —
problème (extrait)}
Soit $f(x) = \dfrac{x^{2}-3x+4}{x-2}$.
\begin{enumerate}[label=\textbf{\alph*)}, leftmargin=*, itemsep=0.15em]
  \item Déterminer l'ensemble de définition de $f$.
  \item Vérifier que $f(x) = x-1+\dfrac{2}{x-2}$.
  \item En déduire que la droite $(D) : y = x-1$ est asymptote oblique à la
        courbe, et étudier la position de la courbe par rapport à $(D)$.
  \item Calculer les limites de $f$ en $2^{-}$ et en $2^{+}$.
\end{enumerate}

\exoann{Baccalauréat, Côte d'Ivoire, série C, session 2021 — exercice 1
(extrait)}
Répondre par vrai ou faux, en justifiant : « toute fonction croissante et
majorée sur $\intervallefo{0}{+\infty}$ admet une limite finie en
$+\infty$ ».
\end{annales}

\begin{corrige}[annales]
\textbf{1. a)} $f(x) = x+\frac{1}{x+1}$ : le second terme tend vers $0$,
donc $f(x) \to -\infty$ en $-\infty$ et $+\infty$ en $+\infty$.
\textbf{b)} En $-1$, le numérateur du second terme vaut $1$ et le
dénominateur tend vers $0$ : négatif à gauche, positif à droite. Donc
$f \to -\infty$ en $-1^{-}$ et $+\infty$ en $-1^{+}$.
\textbf{c)} La droite $x = -1$ est asymptote verticale ; comme
$f(x)-x \to 0$, la droite $y = x$ est asymptote oblique.

\textbf{2. a)} $\Df = \R\setminus\left\{-\frac34\right\}$.
\textbf{b)} $\frac{3x}{3+4x} \to \frac{3x}{4x} = \frac34$ : la droite
$y = \frac34$ est asymptote horizontale en $\pm\infty$.
\textbf{c)} En $-\frac34$, le numérateur tend vers $-\frac94 \neq 0$ et le
dénominateur vers $0$, négatif à gauche et positif à droite :
$f \to +\infty$ en $\left(-\frac34\right)^{-}$ et $-\infty$ en
$\left(-\frac34\right)^{+}$.

\textbf{3. a)} $\Df = \R\setminus\{2\}$.
\textbf{b)} $(x-1)(x-2)+2 = x^{2}-3x+4$. \checkmark
\textbf{c)} $f(x)-(x-1) = \frac{2}{x-2} \to 0$ : la droite $(D)$ est
asymptote oblique. Ce reste est positif pour $x>2$ (courbe au-dessus) et
négatif pour $x<2$ (courbe en dessous).
\textbf{d)} $f \to -\infty$ en $2^{-}$ et $+\infty$ en $2^{+}$.

\textbf{4.} Vrai. Une fonction croissante ne peut que monter ; étant
majorée, elle ne peut pas dépasser son majorant : elle s'accumule donc sous
une valeur limite finie.
\end{corrige}

% ===========================================================================
\begin{vraifaux}
\exo Une courbe ne peut pas couper son asymptote.

\exo Si $\limpinf f(x) = +\infty$, la courbe n'a pas d'asymptote en
$+\infty$.

\exo Une fonction rationnelle a toujours une asymptote verticale.

\exo Si $f(x) = 2x+3+\frac{1}{x^{2}}$, la droite $y = 2x+3$ est asymptote
oblique en $+\infty$.

\exo $\frac{0}{0}$ signifie que la limite vaut $0$.
\end{vraifaux}

\begin{corrige}[vrai ou faux]
\textbf{1.} Faux — une courbe peut couper son asymptote horizontale ou
oblique, parfois plusieurs fois ; elle ne peut simplement pas couper une
asymptote \emph{verticale}, puisque la fonction n'y est pas définie.
\textbf{2.} Faux — elle peut avoir une asymptote \emph{oblique}, ou une
branche parabolique.
\textbf{3.} Faux — $f(x) = \frac{x}{x^{2}+1}$ est définie sur $\R$ tout
entier.
\textbf{4.} Vrai — la différence $\frac{1}{x^{2}}$ tend vers $0$.
\textbf{5.} Faux — c'est une forme indéterminée : la limite peut valoir
n'importe quoi, comme le montrent $\frac{x^{2}}{x} \to 0$ et
$\frac{x}{x^{2}} \to \infty$ en $0$.
\end{corrige}

% ===========================================================================
\begin{jecherche}[une famille de fonctions]
Pour tout réel $m$, on considère la fonction
$f_{m}(x) = \dfrac{mx+1}{x-3}$.

\begin{enumerate}[label=\textbf{\arabic*.}, leftmargin=*, itemsep=0.3em]
  \item Donner l'ensemble de définition de $f_{m}$ ; il ne dépend pas de
        $m$.
  \item Calculer $\limpinf f_{m}(x)$ en fonction de $m$. Quelle asymptote
        cela donne-t-il ?
  \item Existe-t-il une valeur de $m$ pour laquelle la courbe admet l'axe des
        abscisses pour asymptote horizontale ?
  \item Pour $m = 1$, écrire $f_{1}(x)$ sous la forme
        $1+\dfrac{k}{x-3}$ et en déduire les limites en $3^{-}$ et $3^{+}$.
  \item Toutes les courbes de la famille passent-elles par un même point ?
\end{enumerate}
\end{jecherche}

\begin{corrige}[une famille de fonctions]
\textbf{1.} $\Df = \R\setminus\{3\}$ pour toute valeur de $m$.

\textbf{2.} $\limpinf \frac{mx+1}{x-3} = m$ : la droite $y = m$ est
asymptote horizontale.

\textbf{3.} L'axe des abscisses a pour équation $y = 0$ : il faut $m = 0$.
La fonction est alors $f_{0}(x) = \frac{1}{x-3}$.

\textbf{4.} $\frac{x+1}{x-3} = \frac{(x-3)+4}{x-3} = 1+\frac{4}{x-3}$. Le
numérateur $4$ étant positif, on obtient $-\infty$ en $3^{-}$ et $+\infty$
en $3^{+}$.

\textbf{5.} Cherchons un point commun : $f_{m}(x)$ doit être indépendant de
$m$, donc le terme $\frac{mx}{x-3}$ doit disparaître, ce qui exige $x = 0$.
Et alors $f_{m}(0) = \frac{1}{-3} = -\frac13$ pour tout $m$ : toutes les
courbes passent par $A\left(0\,;\,-\frac13\right)$.
\end{corrige}

% ===========================================================================
\begin{jemeteste}
Vingt-cinq minutes.

\begin{enumerate}[label=\textbf{\arabic*.}, leftmargin=*, itemsep=0.3em]
  \item Calculer $\limpinf \dfrac{-3x^{2}+x}{x^{2}+4}$.
  \item Calculer $\displaystyle\lim_{x \to 0^{-}} \dfrac{3}{x}$.
  \item Calculer $\displaystyle\lim_{x \to 4} \dfrac{x^{2}-16}{x-4}$.
  \item Soit $f(x) = \dfrac{2x-5}{x-1}$. Déterminer $\Df$, les deux
        asymptotes, et écrire $f(x)$ sous la forme $2+\frac{k}{x-1}$.
  \item Soit $g(x) = \dfrac{x^{2}+2}{x+1}$. Vérifier que
        $g(x) = x-1+\dfrac{3}{x+1}$, puis donner l'asymptote oblique.
\end{enumerate}
\end{jemeteste}

\begin{corrige}[je me teste]
\textbf{1.} $-3$.
\textbf{2.} $-\infty$.
\textbf{3.} $x+4 \to 8$.
\textbf{4.} $\Df = \R\setminus\{1\}$ ; $x = 1$ et $y = 2$ ;
$\frac{2x-5}{x-1} = \frac{2(x-1)-3}{x-1} = 2-\frac{3}{x-1}$.
\textbf{5.} $(x-1)(x+1)+3 = x^{2}-1+3 = x^{2}+2$ \checkmark ; l'asymptote
oblique est $y = x-1$.
\end{corrige}

% ===========================================================================
\begin{commeaubac}
\bmtsource{Exercice original au format de l'épreuve — série A \textbullet\ 5 points}
Soit $f$ la fonction définie sur $\R\setminus\{1\}$ par
\[
  f(x) = \frac{x^{2}-2x+3}{x-1}.
\]
On note $\Cf$ sa courbe représentative dans un repère orthonormé d'unité
$1$ cm.

\begin{enumerate}[label=\textbf{\arabic*.}, leftmargin=*, itemsep=0.25em]
  \item Justifier que $1$ est la seule valeur interdite. \bareme{0,5 pt}
  \item Vérifier que, pour tout $x \neq 1$,
        $f(x) = x-1+\dfrac{2}{x-1}$. \bareme{1,5 pt}
  \item En déduire que la droite $(D) : y = x-1$ est asymptote à $\Cf$ en
        $+\infty$ et en $-\infty$. \bareme{1 pt}
  \item Étudier la position de $\Cf$ par rapport à $(D)$. \bareme{1 pt}
  \item Calculer les limites de $f$ en $1^{-}$ et en $1^{+}$, et en déduire
        une autre asymptote. \bareme{1 pt}
\end{enumerate}
\end{commeaubac}

\begin{corrige}[comme au baccalauréat]
\textbf{1.} Le dénominateur $x-1$ s'annule pour $x = 1$ et pour cette valeur
seulement.

\textbf{2.} $(x-1)(x-1)+2 = x^{2}-2x+1+2 = x^{2}-2x+3$. \checkmark

\textbf{3.} $f(x)-(x-1) = \dfrac{2}{x-1}$, qui tend vers $0$ en $+\infty$
comme en $-\infty$ : la droite $(D)$ est donc asymptote oblique dans les
deux directions.

\textbf{4.} Le signe de $f(x)-(x-1)$ est celui de $x-1$ : positif pour
$x>1$, négatif pour $x<1$. La courbe est donc au-dessus de $(D)$ sur
$\intervalleo{1}{+\infty}$ et en dessous sur $\intervalleo{-\infty}{1}$.

\textbf{5.} $\frac{2}{x-1} \to -\infty$ en $1^{-}$ et $+\infty$ en $1^{+}$ ;
le terme $x-1$ tend vers $0$ et ne change rien. Donc
$f \to -\infty$ en $1^{-}$ et $+\infty$ en $1^{+}$ : la droite $x = 1$ est
asymptote verticale.
\end{corrige}

\begin{notepourlaclasse}
Le point qui coûte le plus de points à l'examen n'est pas le calcul de la
limite : c'est la phrase d'interprétation qui doit suivre. Faites écrire
systématiquement le couplet complet — « donc la droite d'équation … est
asymptote … à la courbe en … » — même quand l'énoncé ne le demande pas
explicitement ; il le demande presque toujours dans la question suivante.

La forme $f(x) = ax+b+\frac{k}{x-c}$ mérite une séance à elle seule. Elle
revient dans tous les exercices de ce chapitre et dans la moitié de ceux du
chapitre 6 : entraînez la vérification par développement, qui est la
rédaction attendue au baccalauréat, plutôt que la division euclidienne, que
peu d'élèves de la série mènent sans erreur.
\end{notepourlaclasse}
